\documentclass{my_eth_beamer}

\usepackage{my_math_eth}

\numberwithin{equation}{section}

% -----------------------------------------------------------------------------
% this document
\title{Introduction to Survey Sampling}
\newcommand{\TITLE}{Simple Random Sampling}
\newcommand{\NR}{3}
\subtitle{Part \NR: \TITLE}
\author[Tobias Schoch]{Tobias Schoch}
\date{2024}

% -----------------------------------------------------------------------------
\begin{document}

\setstretch{1.1}

\section{\TITLE}\label{sec:srs}
\setcounter{section}{\NR}
\setcounter{Bsp}{10}
\setcounter{Def}{10}
\setcounter{EinSatz}{1}

% -----------------------------------------------------------------------------
\begin{frame}
    \frametitle{Overview}
    \noindent \hspace{0.2em}\textcolor{fhnwblue}{\bfseries \NR~\TITLE}
    \begin{enumerate}[\NR.1]
        \item Estimation
        \item Statistical inference
        \item Sample size determination
        \item Estimation for subpopulations -- Unplanned domains
    \end{enumerate}
\end{frame}

% -----------------------------------------------------------------------------
\begin{frame}
    \frametitle{\insertsectionnumber~\insertsection}
    \begin{itemize}
    \setlength{\itemsep}{10pt}
        \item Simple random sampling (\SRS) is the \alert{gold standard} of
            probability sampling designs.
        \item Let $U$ be a population of $N$ elements. The sample size $n$
            ($0 < n < N$) under \SRS can be expressed by the
            \alert{sampling fraction},
            \begin{equation*}
                f = \frac{n}{N}.
            \end{equation*}
        \item The \alert{sampling design} $p_{\mSRS}$ was introduced in
            Chapter 2 and is given, for $s \in \samplespace$, by
            \begin{equation*}
                p_{\mSRS}(s) = \begin{cases}
                        1 / \binom{N}{n} & \text{if} \; \vert s \vert = n,\\
                        0 & \text{otherwise}.
                    \end{cases}
            \end{equation*}
    \end{itemize}
\end{frame}

% -----------------------------------------------------------------------------
\begin{frame}[fragile]
    \begin{itemize}
    \setlength{\itemsep}{12pt}
        \item The (first-order sample) \alert{inclusion probabilities} under
            $p_{\mSRS}$ have the following form (see Chapter 2):
            \begin{equation*}
                \pi_i = \PP(i \in s) = \frac{n}{N} \qquad \text{for every}
                    \; i \in U.
            \end{equation*}
        \item \alert{Selecting} an \SRS in \textsf{R}:

        \vspace{1em}

\begin{Rcode}
    sample(x, size, replace = FALSE, prob = NULL)
\end{Rcode}

        \item For example,
    \end{itemize}

\begin{Rcode}
    > set.seed(2048)
    > sample(1:100, 10)
    [1] 23 97 88 96 43 27 57 67 31 12
\end{Rcode}

\end{frame}

\subsection{Estimation}

% -----------------------------------------------------------------------------
\begin{frame}
    \frametitle{\insertsectionnumber.\insertsubsectionnumber~\insertsubsection}
    \begin{itemize}
    \setlength{\itemsep}{8pt}
        \item Let $s$ be an \SRS of size $n$ from $U$. The variable of
            interest is denoted by $y_i$.
        \item An \emph{unbiased} estimator of the \alert{population mean}
            $\myu$ is
    \end{itemize}
    \vspace{0.5em}
    \begin{equation*}
        \mys = \frac{1}{n} \sum_{i \in s} y_i.
    \end{equation*}
    \vspace{-0.5em}
    \begin{Exkurs}[Unbiasedness]\label{exkurs:srs_unbiased}
        \vspace{-0.25em}
        To show that $\mys$ is an unbiased estimator of $\myu$, we express
        the estimator as a sum over $i \in U$,
        \begin{equation*}
            \mys = \frac{1}{n} \sum_{i \in s} y_i =
                \frac{1}{n} \sum_{i \in U} \indicator\{i \in s\} \, y_i,
        \end{equation*}
    \end{Exkurs}
\end{frame}

% -----------------------------------------------------------------------------
\begin{frame}
    \begin{Exkurs_forts}[Cont.]
        \vspace{0.25em}
        where $\indicator\{i \in s\}$ is the sample inclusion indicator.
        Under the sampling design, $\indicator\{i \in s\}$ is a Bernoulli
        random variable. Under $p_{\mSRS}$, its expected value is
        \begin{equation*}
            \EE_p\big(\indicator\{i \in s\}\big) = \pi_i = \frac{n}{N}
        \end{equation*}
        for every $i \in U$. Hence,
        \begin{equation*}
            \EE_p(\mys) = \frac{1}{n} \sum_{i \in U}
                \EE_p\big(\indicator\{i \in s\}\big) \, y_i =
                    \frac{1}{n} \sum_{i \in U} \pi_i \, y_i =
                    \frac{1}{n} \sum_{i \in U} \frac{n}{N} \, y_i = \myu,
        \end{equation*}
        \noindent as required.
    \end{Exkurs_forts}
\end{frame}

% -----------------------------------------------------------------------------
\begin{frame}
    \begin{itemize}
    \setlength{\itemsep}{8pt}
        \item The \alert{theoretical variance} of the estimator $\mys$ ($\var$
            or $\var_p$ if we want to highlight the dependence on the design
            $p$) is
            \begin{equation*}
                \var(\mys) = \left(1 - \frac{n}{N}\right)\frac{\varyu}{n} =
                    (1 - f)\frac{\varyu}{n},
            \end{equation*}
            \noindent where $\varyu$ is the population variance of the
            $y$-variable,
            \begin{equation*}
                \varyu = \frac{1}{N - 1}\sum_{i \in U} (y_i - \myu)^2.
            \end{equation*}
        \item The factor $(1 - f)$ is called the \alert{finite population
            correction} (FPC). Without this factor, $\var(\mys)$ would equal
            the familiar variance formula for the sample mean from classical
            statistics.
        \item In the \alert{special case} $n = N$ (a census),
            $\var(\mys) = 0$.
    \end{itemize}
\end{frame}

% -----------------------------------------------------------------------------
\begin{frame}
    \begin{itemize}
    \setlength{\itemsep}{14pt}
        \item An unbiased \alert{variance estimator} for $\mys$ that requires
            only the observed values $y_i$, $i \in s$, is
            \begin{equation*}
                \widehat{\var}(\mys) = \left(1 - \frac{n}{N}\right)
                    \frac{\varys}{n},
            \end{equation*}
            \noindent where
            \begin{equation*}
                \varys = \frac{1}{n - 1}\sum_{i \in s}(y_i - \mys)^2.
            \end{equation*}
        \item The estimated \alert{standard error} $\widehat{\se}(\mys)$ is
            defined as the square root of $\widehat{\var}(\mys)$.
        \item The estimated \alert{coefficient of variation} is, for
            $\mys > 0$,
            \begin{equation*}
                \widehat{\varcoef}(\mys) = \frac{\widehat{\se}(\mys)}{\mys}.
            \end{equation*}
    \end{itemize}
\end{frame}

% -----------------------------------------------------------------------------
\begin{frame}
    \frametitle{\normalsize Population total}
    \begin{itemize}
    \setlength{\itemsep}{8pt}
        \item The \alert{population total} $\ty$ can be estimated unbiasedly by
        \begin{equation*}
            \tys = N \,\mys.
        \end{equation*}
    \end{itemize}

    \vspace{1em}

    \textcolor{fhnwblue}{\textbf{Proportion}}
    \begin{itemize}
    \setlength{\itemsep}{8pt}
        \item Let $b_i = \indicator\{i \in B\}$ be the indicator variable for
            the \alert{dichotomous characteristic} $B$, $i \in U$ (e.g.,
            whether a person is employed, ``yes'' = 1 and ``no'' = 0).
        \item The \alert{population proportion} $\propu_b$ is
            \begin{equation*}
                \propu_b = \frac{1}{N} \sum_{i \in U} b_i =
                    \frac{\text{number of units with characteristic $B$}}{N}.
            \end{equation*}
    \end{itemize}
\end{frame}

\subsection{Statistical inference}\label{sec:srs_inference}

% -----------------------------------------------------------------------------
\begin{frame}
    \frametitle{\insertsectionnumber.\insertsubsectionnumber~\insertsubsection}
    \begin{itemize}
    \setlength{\itemsep}{12pt}
        \item Recall the \alert{central limit theorem} from Chapter 2.
            Under suitable regularity conditions,
            \begin{equation*}
                \frac{\mys - \myu}{\widehat{\se}(\mys)}
            \end{equation*}
            is approximately $N(0,1)$-distributed for sufficiently large $n$.
        \item Thus, for sufficiently large $n$, an approximate
            $100(1 - \alpha)\%$ \alert{confidence interval} (CI) for $\myu$ is
            \begin{equation}\label{eq:srs_ci}
                \mys \; \pm \; z_{1-\alpha/2}\,
                    \widehat{\se}(\mys),
            \end{equation}
            where $z_{1-\alpha/2}$ is the $(1-\alpha/2)$ quantile of the
            $N(0,1)$ distribution.
        \end{itemize}
\end{frame}

% -----------------------------------------------------------------------------
\begin{frame}
    \begin{itemize}
    \setlength{\itemsep}{8pt}
        \item The confidence interval is based on a \alert{normal
            approximation}.
        \vspace{0.25em}
            \begin{itemize}
            \setlength{\itemsep}{6pt}
                \item There is no \emph{universal} minimum sample size for this
                    approximation to be accurate. Its accuracy depends on the
                    finite population, in particular on skewness and extreme
                    values.
                \item Because such features are \emph{common} in survey
                    applications, simple rules such as $n \geq 30$'' or $n \geq
                    50$'' should be used with caution.
            \end{itemize}
        \item A \alert{practical question} is therefore how large $n$ should be
            for the normal approximation to be sufficiently accurate.
        \item For populations with marked \alert{positive skewness},
            \citet[][p.~42]{cochran_1977} proposed the \alert{crude rule}
            \begin{equation*}
                n > 25,G_1^2,
            \end{equation*}
            where $G_1$ denotes the population skewness.
    \end{itemize}
\end{frame}

% -----------------------------------------------------------------------------
\begin{frame}
    \begin{itemize}
    \setlength{\itemsep}{12pt}
        \item Population skewness is defined as
            \begin{equation*}
                G_1 = \frac{1}{N} \sum_{i \in U} \left(\frac{y_i -
                \myu}{\sigma_y}\right)^3
            \end{equation*}
        \item Cochran emphasized that this is only a \alert{crude rule of
            thumb}, not a general criterion for the normal approximation.
        \item \citet{sugden_smith_etal_2000} proposed the more conservative rule
            \begin{equation*}
                n_{\min} = 28 + 25\,G_1^2.
            \end{equation*}
    \end{itemize}
\end{frame}

% -----------------------------------------------------------------------------
\begin{frame}
    \frametitle{\normalsize Confidence interval for an estimated proportion}
    \begin{itemize}
    \setlength{\itemsep}{12pt}
        \item To construct a $100(1 - \alpha)\%$ CI for the estimated
            proportion $\props$, we can replace $\mys$ by $\props$ in
            (\ref{eq:srs_ci}), giving
            \begin{equation*}
                \props \, \pm \, z_{1 - \alpha / 2} \, \sqrt{(1 -f) \,
                    \frac{\props (1 - \props)}{n - 1}}.
            \end{equation*}
        \item The resulting CI is called the \alert{Wald interval}. It can be
            used under the usual conditions, i.e., when $n \propu \geq 5$
            and $n(1 - \propu) \geq 5$.
        \item The Wald interval can behave erratically and is considered
            unreliable when $n$ is small or $\propu$ is close to $0$ or $1$.
            \citet{brown_cai_etal_2001} therefore recommend the
            \alert{Wilson interval}; it is implemented in the function
            \code{svyciprop()} of the \textsf{R} package \code{survey}.
    \end{itemize}
\end{frame}

\subsection{Sample size determination}

% -----------------------------------------------------------------------------
\begin{frame}
    \frametitle{\insertsectionnumber.\insertsubsectionnumber~\insertsubsection}
    \begin{itemize}
    \setlength{\itemsep}{12pt}
        \item One of the first questions in survey planning is:
            \alert{How large should the sample be?}
        \item The answer depends mainly on three \alert{factors}:
            \begin{itemize}
            \setlength{\itemsep}{6pt}
                \item the total survey budget;
                \item the cost of observing one unit;
                \item the required precision of the estimates.
            \end{itemize}
        \item In this section, we consider only \alert{sample size
            determination} based on the \alert{required precision}.
    \end{itemize}
\end{frame}

% -----------------------------------------------------------------------------
\begin{frame}
    \begin{itemize}
    \setlength{\itemsep}{10pt}
        \item Even when a survey collects \alert{many variables}, it is often
            possible to focus on one or two variables of \alert{primary
            interest} and use them to determine the sample size.
        \item In what follows, we assume that the required precision is
            specified for the estimated \alert{population mean} $\myu$.
        \item Let $\epsilon > 0$ denote the \alert{margin of error}. The
            precision requirement can then be stated as follows:
            \alert{Choose the sample size such that}
            \begin{equation}\label{eq:srs_abs_err}
                \PP\big(\vert \mys - \myu \vert > \epsilon \big) \leq \alpha,
            \end{equation}

            \vspace{-0.25em}

            \noindent where the probability $\alpha$ (e.g., 5\%) and
            $\epsilon > 0$ are specified in advance.
        \item By the \alert{central limit theorem} (see Theorem 1 in Chapter
            2), $(\mys - \myu) \, / \, \widehat{\se}(\mys)$ has an
            approximately $N(0,1)$ distribution for a sufficiently large
            sample size.
    \end{itemize}
\end{frame}

% -----------------------------------------------------------------------------
\begin{frame}
    \begin{itemize}
    \setlength{\itemsep}{10pt}
        \item Using the $N(0,1)$ approximation, we can rewrite
            (\ref{eq:srs_abs_err}) and obtain the \alert{initial estimate}
            that the sample size must be at least
            \begin{equation}\label{eq:srs_n}
                n_0 = \left(\frac{z_{1 - \frac{\alpha}{2}} \cdot
                    \sigma_y}{\epsilon}\right)^2,
            \end{equation}
            \noindent where $z_{1 - \alpha/2}$ is the $(1 - \alpha/2)$
            quantile of the $N(0,1)$ distribution and $\sigma_y$ is the
            standard deviation of the $y$-variable.
        \item \alert{Alternative derivation}: Set the half-width of the
            $100(1 - \alpha)\%$ confidence interval for $\mys$ equal to
            $\epsilon$:
                \begin{equation*}
                    \mys \, \pm \, \underbrace{z_{1 - \alpha / 2} \;
                        \widehat{\se}(\mys)}_{=\epsilon}.
                \end{equation*}
        \item In practice, $n_0$ is \alert{rounded up} to the next integer.
    \end{itemize}
\end{frame}

% -----------------------------------------------------------------------------
\begin{frame}
    \begin{itemize}
    \setlength{\itemsep}{12pt}
        \item If $n_0 > N$, we set $n_0 = N$ (a census).
        \item The estimate in (\ref{eq:srs_n}) is suitable when the resulting
            sampling fraction $n_0/N$ is negligible (e.g., $n_0/N < 0.01$).
            Otherwise, we use the \alert{refined estimate}
            \begin{equation}\label{eq:srs_n_update}
                n = \frac{n_0}{1 + n_0 / N},
            \end{equation}

            \vspace{0.5em}

            \noindent which accounts for the \alert{finite population
            correction}.
        \item If the population is very large, typically $n \ll N$.
            \vspace{0.25em}
            \begin{itemize}
            \setlength{\itemsep}{6pt}
                \item[\goto] In this case, the required sample size is
                    determined by the \alert{absolute size} of $n$, not by the
                    ratio $n/N$.
            \end{itemize}

    \end{itemize}
\end{frame}

% -----------------------------------------------------------------------------
\begin{frame}
    \begin{itemize}
    \setlength{\itemsep}{8pt}
        \item In practice, the (population) standard deviation $\sigma_y$ in
            (\ref{eq:srs_n}) is \alert{unknown} and must be estimated. The
            following options are available:
            \vspace{0.5em}
            \begin{enumerate}
            \setlength{\itemsep}{6pt}
                \item conduct a \alert{pilot study};
                \item use information from \alert{previous studies} or the
                    literature;
                \item use a \alert{rough estimate} if the first two options
                    are not feasible.
            \end{enumerate}
    \end{itemize}
\end{frame}

% -----------------------------------------------------------------------------
\begin{frame}
    \frametitle{\normalsize Anticipated nonresponse}
    \begin{itemize}
    \setlength{\itemsep}{12pt}
        \item \alert{Nonresponse and refusals} are common in survey practice.
        \item The \alert{expected amount} of nonresponse must therefore be
            taken into account at the planning stage.
        \item Let $r$ denote the anticipated response rate, $r \in (0,1]$.
            The \alert{gross sample size} $n^*$ is
            \begin{equation*}
                n^* = \frac{n}{r},
            \end{equation*}
            \noindent where $n$ is the \alert{net sample size} determined
            using the formulas above.
        \item In practice, the response rate $r$ can usually only be estimated
            \alert{roughly} from previous experience.
    \end{itemize}
\end{frame}

\subsection{Estimation for subpopulations -- Unplanned domains}

% -----------------------------------------------------------------------------
\begin{frame}
    \frametitle{\insertsectionnumber.\insertsubsectionnumber~\insertsubsection}
    \begin{itemize}
    \setlength{\itemsep}{8pt}
        \item Let $U_d \subset U$ be a \alert{subpopulation} of $U$ that is of
            particular interest for the analysis (a domain of interest).
        \item We also refer to $U_d$ simply as a \alert{domain}.
    \end{itemize}
    \vspace{2em}
    \begin{columns}
        \begin{column}{0.5\textwidth}
            \begin{tabular}{lrr}
                \toprule
                \addlinespace
                Language region & $\quad$ Number & $\quad$ Proportion\\
                \addlinespace
                \midrule
                German     & 1\,260 &  66\%\\
                French     &    572 &  30\%\\
                Italian    &     70 &   4\%\\
                \midrule
                Total      & 1\,902 & 100\%\\
                \bottomrule
            \end{tabular}
        \end{column}
        \begin{column}{0.3\textwidth}

            \vspace{-5.5em}

            {\small
                \textcolor{fhnwblue}{\textbf{Example.}}
                Population of ophthalmologists in Switzerland with a license
                to practice, by language region
                (MedReg, 2022)
            }
        \end{column}
    \end{columns}
\end{frame}

% -----------------------------------------------------------------------------
\begin{frame}
    \begin{itemize}
    \setlength{\itemsep}{4pt}
        \item Assume that domain $U_d$ is \alert{unplanned}. This means that
            \vspace{0.25em}
            \begin{itemize}
                \setlength{\itemsep}{6pt}
                \item the elements of $U_d$ \alert{cannot be identified} in
                    advance, and
                \item its population size $N_d$ is \alert{unknown}
                    ($P_d = N_d/N$ is unknown).
            \end{itemize}
        \item We want to estimate the \alert{mean} $\widebar{y}_{U_d}$ in the
            \alert{subpopulation} $U_d$.
        \item First, note that $\widebar{y}_{U_d}$ can be written as
            \begin{equation*}
                \widebar{y}_{U_d} =
                    \frac{\sum_{i \in U} \indicator_{id} \; y_i}{\sum_{i \in U}
                    \indicator_{id}},
            \end{equation*}

            \vspace{-0.5em}

            \noindent where
            \begin{equation*}
                \indicator_{id} = \indicator\{i \in U_d\} = \begin{cases}
                        1 & \text{if} \; i \in U_d,\\
                        0 & \text{otherwise}.
                    \end{cases}
            \end{equation*}
            \noindent Thus, $\widebar{y}_{U_d}$ is the \alert{ratio} of two
            population totals.
    \end{itemize}
\end{frame}

% -----------------------------------------------------------------------------
\begin{frame}
    \begin{itemize}
    \setlength{\itemsep}{8pt}
        \item The \alert{sample size} $n_d$ of the subsample $s_d$ falling in
            $U_d$ ($s_d = s \cap U_d$) is known \alert{only after} the sample
            $s$ has been selected.
        \item Using the indicator variable
            \begin{equation*}
                \indicator_{id}^s = \indicator\{i \in s_d\} = \begin{cases}
                    1 & \text{if} \; i \in s_d,\\
                    0 & \text{otherwise},
                \end{cases} \qquad i \in U,
            \end{equation*}
            \noindent the \alert{domain sample size} can be written as
            \begin{equation*}
                n_d = \sum_{i \in U}\indicator_{id}^s.
            \end{equation*}
        \item Since $\indicator_{id}^s$ is a function of the random sample
            $s$, $n_d$ is a \alert{random variable} with expected value
            $nP_d$ (hypergeometric distribution).
    \end{itemize}
\end{frame}

% -----------------------------------------------------------------------------
\begin{frame}
    \begin{Bsp}[Ophthalmologists]\label{ex:srs_augenarzt}
        \vspace{-0.5em}
        \noindent Suppose an \SRS of size $n = 500$ is drawn from the
        population of ophthalmologists. The expected realized sample size
        $n_d$ for Italian-speaking Switzerland is approximately $18$. With
        probability 95\%, $n_d$ lies in the interval $[12,\,26]$.
    \end{Bsp}
    \vspace{0.5em}
    \begin{itemize}
    \setlength{\itemsep}{8pt}
        \item An estimator of $\widebar{y}_{U_d}$ is
            \begin{equation*}
                \widebar{y}_{d}
                =
                \frac{\sum_{i \in s_d} y_i}{n_d},
                \qquad n_d > 0.
            \end{equation*}
        \item This estimator is a \alert{ratio} of two estimated totals; see
            Chapter 7.
        \item Although this estimator has the form of a ratio of two estimated
            totals, it has a useful special property under \SRS.
        \item \alert{Conditional} on $n_d = m > 0$, the domain sample $s_d$ is
            an \alert{\SRS of size $m$ from $U_d$}.
    \end{itemize}
\end{frame}

% -----------------------------------------------------------------------------
\begin{frame}
    \begin{itemize}
    \setlength{\itemsep}{12pt}
        \item Hence,
            \begin{equation*}
                \EE_p\big(\widebar{y}_d \mid n_d = m\big) =
                    \widebar{y}_{U_d}.
            \end{equation*}
        \item Thus, $\widebar{y}_d$ is unbiased for the domain mean,
            \alert{conditional} on observing at least one unit from the domain.
        \item The main difficulty with an \alert{unplanned domain} is that
            $n_d$ is random.
            \vspace{0.25em}
            \begin{itemize}
            \setlength{\itemsep}{8pt}
                \item If the domain is small, $n_d$ may be small, resulting in
                    \alert{low precision}.
                \item It is even possible that $n_d = 0$, in which case
                    $\widebar{y}_d$ cannot be computed.
            \end{itemize}
        \item If reliable estimates are required for a particular domain, it
            should, whenever possible, be taken into account when designing the
            sample (\alert{planned domain}; see Chapter 4).
    \end{itemize}
\end{frame}

% -----------------------------------------------------------------------------
\begin{frame}
    \frametitle{\normalsize Appendix: Determining the required sample size
    (derivation)}
     \begin{itemize}
     \setlength{\itemsep}{6pt}

        \item Let $\props$ denote the estimator of the population proportion
            $\propu$ under the design $p_{\mSRS}$. Ignoring the finite
             population correction, its estimated standard error is
            \begin{equation*}
                \widehat{\se}\big(\props,\big) =
                \sqrt{\frac{\props(1-\props)}{n_0-1}}
                \approx
                \sqrt{\frac{\props(1-\props)}{n_0}}, \tag{C.1}
            \end{equation*}
            where the approximation is appropriate when $n_0$ is sufficiently
             large.

        \item The precision requirement can now be stated as follows:
            \alert{Choose the sample size such that}
            \begin{equation*}
                \PP\big(\vert \props - \propu \, \vert > \epsilon \big) \leq
                \alpha, \tag{C.2}
            \end{equation*}

            \vspace{-0.25em}

            \noindent where the probability $\alpha$ (e.g., 5\%) and the
             \alert{margin of error} $\epsilon > 0$ are specified in advance.
     \end{itemize}
\end{frame}

% -----------------------------------------------------------------------------
\begin{frame}
    \begin{itemize}
    \setlength{\itemsep}{10pt}
        \item The requirement in (C.2) is \alert{equivalent} to
            \begin{equation*}
                \PP\bigl(\lvert\props-\propu\rvert\leq\epsilon\bigr)
                \geq 1-\alpha.
            \end{equation*}
        \item For sufficiently large sample sizes, the \alert{central limit
            theorem} gives
            \begin{equation*}
                Z = \frac{\props-\propu}
                {\widehat{\se}(\props)}
                \approx N(0,1).
            \end{equation*}
        \item Hence, using the \alert{symmetry} of the standard normal
            distribution,
            \begin{equation*}
                \frac{\epsilon}{\widehat{\se}(\props)}
                \geq z_{1-\alpha/2},
                \tag{C.3}
            \end{equation*}
            where $z_{1-\alpha/2}$ is the $(1-\alpha/2)$ quantile of the $N(0,1)$
            distribution.
    \end{itemize}
\end{frame}

% -----------------------------------------------------------------------------
\begin{frame}
    \begin{itemize}
    \setlength{\itemsep}{0pt}
        \item Substituting the approximation from (C.1) into (C.3) gives the
            initial sample size
            \begin{equation*}
                n_0 =
                \frac{P_0(1-P_0)\,z_{1-\alpha/2}^2}{\epsilon^2}, \tag{C.4}
            \end{equation*}
            where $P_0$ is an \alert{anticipated value} of the population
            proportion. If the sampling fraction $n_0/N$ is not negligible, the
            \alert{finite population correction} gives
            \begin{equation*}
                n =
                \frac{n_0}{1+(n_0-1)/N}. \tag{C.5}
            \end{equation*}
        \item Substituting (C.4) into (C.5) gives
            \begin{equation*}
                n \geq \frac{N P_0(1-P_0)z_{1-\alpha/2}^2}{
                    \epsilon^2(N-1) + P_0(1-P_0)z_{1-\alpha/2}^2}.
            \end{equation*}
        \item The resulting sample size is rounded up to the next integer.
    \end{itemize}
\end{frame}

% -----------------------------------------------------------------------------
\setcounter{section}{\the\numexpr\value{section}-1\relax}
\begin{frame}
    \frametitle{Literature}
    {\footnotesize
        \setstretch{1.1}
        \bibliographystyle{fhnw_en}
        \bibliography{master}
    }
\end{frame}
% -----------------------------------------------------------------------------
\lastpage

\end{document}
